
\section{Cubesets}
In this section, we introduce the category of cubes and cubesets, study their basic properties and compare them to $\bbGamma$-sets.



\subsection{The Building Blocks}\label{ssec:index-cats}

In this subsection we first want to discuss several variants of the indexing category for cubesets and relations between candidates of such.

\subsubsection{Defining the Indexing Category}

Let us start with the classical category of combinatorial cubes used in nilspace theory, being introduced e.g. in \cite[p. 3]{Camarena2010}, \cite[Definition 3.1]{Gutman2020} and \cite[Definition 1.1.1]{Candela2017}.

\begin{definition}\label{def:cube_cat}
Define the \emph{cube category $\bbox$} as the skeletal category%
\footnote{If not stated otherwise, by category we mean a (static/ordinary) $(1,1)$-category.}
whose objects are the sets $\{0,1\}^n$ for $n\in\N_0$ and where a morphism $f\colon\{0,1\}^n\to\{0,1\}^m$
is a set-theoretical map that is the restriction of an affine linear map $\Z^n\to\Z^m$.
\end{definition}

To understand this category it is helpful to first study an important subcategory.

\begin{definition}\label{def:linear_cubes}
Denote by $\bblox$ the wide subcategory of $\bbox$ consisting of all maps that are restrictions of linear maps.
We denote this corresponding inclusion by $j\colon\bblox\to\bbox$.
\end{definition}

This linear cube category relates very nicely to a classically known indexing category from (higher) category theory.

\begin{definition}\label{def:indexcats}
Let us define the following small categories.
\begin{enumerate}[label=(\roman*), ref=\thedefinition(\roman*)]
    \item\label{def:finstar}
        Define \emph{Segal's category} $\bbGamma$ as the opposite of $\Fin_\ast$.
        The category $\Fin_\ast$ is (the skeleton of) the category of finite pointed sets.
        This means that its objects are indexed by the set $\N_0$, which we also denote by
        \[\langle n \rangle=\{\ast\}\sqcup\langle n\rangle^\circ=\{\ast\}\sqcup \{1,\ldots,n\}\]
        for $n\in \N_0$.
        Morphisms $\langle n\rangle \to \langle m\rangle$ in $\Fin_*$ are given by base-point preserving maps of sets,
        or equivalently, by partially defined maps $\langle n\rangle^\circ\to \langle m\rangle^\circ$.
    \item\label{def:inert_active}
        We call a morphism $\alpha$ in $\Fin_\ast$ \emph{inert} (or \emph{lower face}) if for all $i\neq\ast$ the preimage $\alpha^{-1}(i)$ consists of exactly one element.
        And we call $\alpha$ \emph{active} (or \emph{diagonal}) if it is everywhere defined, i.e., $\alpha(j)\neq\ast$ for all $j\ne \ast$.
    \item
        Define the category $\finert$ to be the wide subcategory of $\Fin_*$ consisting of inert morphisms only.
        Define $\inert$ as the opposite of $\finert$.
    \item\label{def:almost-inert}
        Define the category $\Fin_{*,\le 1}$ to be the subcategory of $\Fin_*$ where morphisms are required to have fibers of size at most one for all non-basepoints.    
\end{enumerate}
\end{definition}


\begin{proposition}\label{lem:iso-gamma-bblox}
There is a canonical equivalence of categories $\bbGamma\simeq{\bblox}$.
\end{proposition}

To prove this, let us first describe the morphisms of $\bbox$ in a more coordinate-dependent version.

\begin{lemma}\label{lem:coordinate_description}
A set-theoretical map $f\colon\{0,1\}^n\to\{0,1\}^m$ is a morphism in $\bbox$
if and only if it is of the form (its \emph{canonical representation})
\[
    f(x_1,\ldots,x_n) = (y_1,\ldots,y_m)
\]
where each $y_i$ is either $0$, $1$, or there exists a unique $j(i)$ such that $y_i = x_{j(i)}$ or $y_i = 1 - x_{j(i)}$.

Similarly, a morphism $f\colon\{0,1\}^n\to\{0,1\}^m$ in $\bblox$ can be described as
\[
    f(x_1,\ldots,x_n) = (y_1,\ldots,y_m)
\]
where each $y_i$ is either equal to zero or there exists a unique $j(i)$ such that $y_i = x_{j(i)}$.
\end{lemma}

\begin{proof}
The first part can be found in \cite[Lemma 1.1.3]{Candela2017} and \cite[Remark after Definition 3.1]{Gutman2020}.
Let us recall the argument for completeness.

First, if $f$ is of the form in the lemma, then by setting $g_i(x) = f_i(x) - f_i(0)$
we see that $g_i(x)$ is either the zero map, or $g_i(x) = \pm x_{j(i)}$ for some $j(i)$.
In particular, $f$ is the restriction of an affine map from $\Z^n$ to $\Z^m$.

On the other hand, if $f$ arises from an affine map,
then there exists a matrix $A\in\Z^{m\times n}$ and a point $p\in\Z^m$ such that
\[
    f(x) = Ax + p.
\]
From the condition $f(x)\in\{0,1\}^m$, hence $f_i(x)\in\{0,1\}$, for all $x\in\{0,1\}^n$ and all $i\in\{1,\ldots,m\}$ we know that
every row of $A$ has at most one non-zero entry which is either $1$ or $-1$.
Clearly each non-zero entry of $A$ can only be $1$ or $-1$ or else $f_i(x)$ wouldn't land in $\{0,1\}$.
Moreover, if there were two or more non-zero entries, then testing with appropriate $x$ would contradict $f_i(x)\in\{0,1\}$.
Now if in row $i$ every entry is zero, $p_i\in\{0,1\}$ and we set $y_i = p_i$.
If in row $i$ the entry in column $j(i)$ is $1$, $p_i$ has to be zero and we set $y_i = x_{j(i)}$.
If in row $i$ the entry in column $j(i)$ is $-1$, $p_i$ has to be one and we set $y_i = 1 - x_{j(i)}$.
Here, $j(i)$ is uniquely determined.
One can now easily check that the map
\[
    (x_1,\ldots,x_n) \mapsto (y_1,\ldots,y_m)
\]
is the same as the map $f$.

The second part follows from the first noting that an affine map $A$ is linear if and only if $A(0)=0$.
\end{proof}

With this at hand, we can prove that the linear part of $\bbox$ simply is $\bbGamma$.

\begin{proof}[Proof of Proposition \ref{lem:iso-gamma-bblox}]
There is a functor
\[
    F\colon{\bblox}\to\Fin_\ast^\op
\]
which acts on objects as $F(\{0,1\}^n) = \langle n\rangle$.
To a morphism $f\colon\{0,1\}^n\to\{0,1\}^m$ it assigns the morphism
\[
    F(f)\colon\langle m\rangle\to\langle n\rangle,\quad i\mapsto
        \begin{cases}
            j(i),\quad&\text{if}\; y_i = x_{j(i)}, \\
            \ast,\quad&\text{if}\; y_i = 0
        \end{cases}
\]
in $\Fin_\ast$.
In the other direction, there exists the functor
\[
    G\colon\Fin_\ast^\op\to{\bblox}
\]
which acts on objects as $G(\langle n\rangle) = \{0,1\}^n$.
To a morphism $\alpha\colon\langle m\rangle\to\langle n\rangle$ in $\Fin_\ast$ it assigns the morphism
\[
    G(\alpha)\colon\{0,1\}^n\to\{0,1\}^m,\quad (x_1,\ldots,x_n)\mapsto (y_1,\ldots,y_m)
\]
where $y_i = 0$ if $\alpha(i) = \ast$ and $y_i = x_{\alpha(i)}$ otherwise.%
\footnote{Note that this is simply pullback with $\alpha$, if one uses the convention $x_\ast=0$.}
It is a straightforward computation to show that $F$ and $G$ are inverse to each other.
\end{proof}



To reach a categorical description of the whole cube category $\bbox$,
the idea is to note that in fact $\bblox$ is not just a subcategory of $\bbox$ but a retract,
i.e., there is a functor $\bbox\to \bblox$ assigning to every affine linear map its \emph{linear part}.
This exhibits the category $\bbox$ as \emph{fibered over $\bbGamma=\bblox$},
where the fiber of $\langle n\rangle$ consists of the additional \emph{flips} induced by the non-linear base-point translation.
This leads to the description of $\bbox$ as a certain \emph{skew-product} of categories,
or, in more modern language, as \emph{unstraightening} via the classical \emph{Grothendieck construction}.


\begin{construction}\label{cons:unstraightening}
Let us recall the Grothendieck construction.
Given a (pseudo-)functor $F\colon\cC\to\Cat$ one can construct a new category
\[
    \mathrm{Un}(F)=\int_{X\in\cC} F(X)
\]
as follows:
\begin{itemize}
    \item its objects are pairs $(X,A)$ where $X\in\cC$ and $A\in F(X)$ and
    \item a morphism $(X,A)\to (Y,B)$ is a pair
    \[(f\colon X\to Y,\phi\colon F(f)(A)\to B)\]
where $f$ is a morphism in $\cC$ and $\phi$ is a morphism in $F(Y)$.
    \end{itemize}
Now composition is given in the only way possible:
given $(f,\phi)\colon (X,A)\to (Y,B)$ and $(g,\psi)\colon (Y,B)\to (Z,C)$ its composite is the pair
\[
    (g\circ f\colon X\to Y\to Z,\psi\circ F(g)(\phi)\colon F(g\circ f)(A)\to F(g)(B)\to C).
\]
This construction gives rise to a ($\infty$-)functor
\[
    \int_\cC\colon\Fun(\cC,\Cat)\to\Cat/\cC
\]
which induces an equivalence between $\Fun(\cC,\Cat)$ and the subcategory of $\Cat/\cC$
of Grothendieck opfibrations (also called cocartesian fibrations), called the \emph{straightening-unstraightening equivalence}, see \cite[Chapter 3]{Lurie2009}.
\end{construction}


Hence, let us consider an appropriate functor $\bbGamma\to \Cat$ encoding the flip-action of cubes.

\begin{definition}
The assignment $\langle n\rangle\mapsto\F_2^n$ defines a functor $\Fin_\ast^\op\to\Grp$%
\footnote{Note that we denote the group $\Z/2\Z$ by $\mathbb{F}_2$
since we want to avoid the double quotient notation $\ast/(\Z/2\Z)$.}
since a morphism $\alpha\colon\langle n\rangle\to\langle m\rangle$ in $\bbGamma$ induces
by pullback at its opposite a group homomorphism $\alpha_\ast=(\alpha^\op)^\ast\colon\F_2^n\to\F_2^m$.
At coordinate $j\in\{1,\ldots,m\}$ it is given by
\[
    (x_1,\ldots,x_n)\mapsto\begin{cases}
                                x_{\alpha(j)},\quad&\text{if}\;\alpha(j)\neq\ast, \\
                                0,\quad&\text{else}.
                            \end{cases}
\]
By composing with the functor $\Grp\to\Cat,\;G\mapsto */G=\mathrm{B}G$ we obtain a functor
\[
    \ast/\F_2^\bullet\colon\bbGamma\to\Cat,\quad\langle n\rangle\mapsto\ast/\F_2^n.
\]
\end{definition}

With this definition we arrive at the following categorical definition of the cube category.

\begin{proposition}\label{prop:cube_cat_linear}
The cube category $\bbox\to\bbGamma$ is the unstraightening
\[
    \mathrm{Un}(\ast/\F_2^\bullet)=\int_{\bbGamma} \ast/\F_2^\bullet\to\bbGamma
\]
of the functor $\ast/\F_2^\bullet\colon\bbGamma\to\Cat$.
\end{proposition}

In our case, the Grothendieck construction takes the following form.

\begin{remark}
Let us describe the category
\[
    \int_{\bbGamma} \ast/\F_2^\bullet\to\bbGamma.
\]
Since the category $\ast/\F_2^n$ consists of exactly one object $\ast_n$,
an object in ${\int_{\bbGamma}} \ast/\F_2^\bullet$ is just an object $\langle n\rangle$ of $\bbGamma$.

A morphism $(f,\phi)\colon\langle n\rangle\to\langle m\rangle$ consists of a morphism $f$ in $\bbGamma$
and a morphism $\phi\colon(\ast/\F_2^\bullet)(f)(\ast_n)=\ast_m\to\ast_m$ in $\ast/\F_2^m$ which corresponds to an element $a\in\F_2^m$.

Now composition works as follows: given $(f,a)\colon\langle n\rangle\to\langle m\rangle$
and $(g,b)\colon\langle m\rangle\to\langle k\rangle$, in the first component it is just ordinary composition $g\circ f$.
In the second component, the element of $\F_2^k$ is given by $b+g_\ast(a)$, so
\[
    (g,b)\circ(f,a)=(g\circ f, b+g(a)).
\]
\end{remark}

Note that the functor $\ast/\F_2^\bullet\colon\bbGamma\to\Cat$ factors over the isomorphism
$G\colon\bbGamma\to\bblox$ from Proposition~\ref{lem:iso-gamma-bblox} by interpreting $\{0,1\}^n$
as the group $\F_2^n$.

\begin{remark}
Since $\Aut_{\bbGamma}(\langle n\rangle) = \uS_n$ and $\Aut_{\ast/\F_2^n}(\ast) = \F_2^n$
there is a split short exact sequence
\begin{center}
\begin{tikzcd}
	0 & {\F_2^n} & {\Aut_{\int_{\bbGamma} \ast/\F_2^\bullet}(\langle n\rangle)} & {\uS_n} & 0
	\arrow[from=1-1, to=1-2]
	\arrow[from=1-2, to=1-3]
	\arrow[from=1-3, to=1-4]
	\arrow[from=1-4, to=1-5]
\end{tikzcd}
\end{center}
which exhibits $\Aut_{\int_{\bbGamma} \ast/\F_2^\bullet}(\langle n\rangle)$ as the group-theoretic semidirect product of $\F_2^n$ and $\uS_n$.
\end{remark}


The group $\F_2^n$ now acts as \emph{base-point shift} or, rather, as \emph{flip} on the $n$-cubes.
\begin{lemma}\label{lem:linearization-auto-bblox}
For every $n\in\N_0$ there exists an injective group homomorphism
\[
    \phi_{(-)}\colon\F_2^n\to\Aut_{{\bbox}}(\{0,1\}^n)
\]
with $\phi_a(0) = a$.
This group homomorphism is natural in ${\bblox}$ in the sense that for every map
$f\colon\{0,1\}^n\to\{0,1\}^m$ in ${\bblox}$ we have that $\phi_{f(a)}f = f\phi_a$ for all $a\in\{0,1\}^n$.

Moreover, for every morphism $f\colon\{0,1\}^n\to\{0,1\}^m$ there exists a unique factorization
\begin{center}
\begin{tikzcd}
	{\{0,1\}^n} & {\{0,1\}^m} & {\{0,1\}^m}
	\arrow["g", from=1-1, to=1-2]
	\arrow["\phi", from=1-2, to=1-3]
\end{tikzcd}
\end{center}
where $g$ is a morphism in ${\bblox}$ and $\phi$ is in the image of $\phi_{(-)}$.
In this case, $\phi = \phi_{f(0)}$.
\end{lemma}

\begin{proof}
Let $a\in\{0,1\}^n$.
Define $\phi_{a,i}(x_1,\ldots,x_n) = x_i$ if $a_i = 0$ and $\phi_{a,i}(x_1,\ldots,x_n) = 1-x_i$ if $a_i = 1$.
Thus, identifying $\{0,1\}^m$ with $\F_2^m$ the map $\phi_a$ simply is addition of $a$.
Hence, $\phi_a$ is self-inverse and thus an automorphism of $\{0,1\}^n$.
Moreover, $\phi_a\phi_b = \phi_{a+b}$ where the sum $a+b$ is taken in $\F_2^n$:
If $\phi_a$ is the identity, then $a = \phi_a(0) = 0$ by definition, so $\phi_{(-)}$ is injective.


To prove the second part, let $f\colon\{0,1\}^n\to\{0,1\}^m$ be a map in $\bblox$.
We need to show $\phi_{f(a)}f\phi_a = f$.
The idea is to compare coordinatewise -- either $f_i$ is zero and then $\phi_{f(a)}$ does not affect the $i$-th coordinate, resulting in $f$,
or both $\phi_a$ and $\phi_{f(a)}$ do the same thing (do or don't contribute a flip),
again resulting in $f$.

For $1\leq i\leq m$ we have that $f_i(x_1,\ldots,x_n) = y_i$ where $y_i \equiv 0$ or $y_i = x_{j(i)}$ for some unique $j(i)$ by Lemma~\ref{lem:coordinate_description}.
In the first case $f_i(x) = 0$, we have $f_i(a) = 0$, so $\phi_{f(a),i}(y_1,\ldots,y_m) = y_i$, hence
\[
    \phi_{f(a),i}(f(\phi_a(x_1,\ldots,x_n)))=f_i(\phi_a(x_1,\ldots,x_n)) = 0=f_i(a).
\]
In the second case $f_i(x) = x_{j(i)}$.
If $a_{j(i)}=0$, then $f_i(a)=0$ and so $\phi_{a,j(i)}(x)=x_{j(i)}$,
$f_i(x)= x_{j(i)}$ and $\phi_{f(a),i}(y) = y_i$, therefore
\[
    \phi_{f(a),i}(f(\phi_a(x))) = x_{j(i)} = f_i(x).
\]
If $a_{j(i)}=1$, then $f_i(a)=1$ and so $\phi_{a, j(i)}(x)=1-x_{j(i)}$,
$f_i(x)=x_{j(i)}$ and $\phi_{f(a),i}(y) = 1-y_i$, therefore
\[
    (\phi_{f(a)}(f(\phi_a))(x))_i = \phi_{f(a),i}(1-x_{j(i)}) = x_{j(i)} = f_i(x).
\]
This shows that indeed $\phi_{f(a)}f\phi_a = f$.

For the last part let $g = \phi_{f(0)}f$ which is linear since $g(0) = \phi_{f(0)}(f(0)) = 0$.
By construction, $\phi_{f(0)}$ is self-inverse and so $f$ factors as $\phi_{f(0)}\circ g$.
For uniqueness, take another factorization of $f = \phi'g'$ where $g'$ is linear and $\phi'$ is in the image of $\phi_{(-)}$.
Since $g'$ is linear, $g'(0)=0$ and we obtain that $\phi'(0) = f(0) = \phi_{f(0)}(0)$.
Therefore, $\phi'=\phi_{f(0)}$ and from self-inverseness of $\phi$ it follows that $g' = g$.
\end{proof}

\begin{proof}[Proof of Proposition~\ref{prop:cube_cat_linear}]
We obtain a functor
\[
    H\colon{\bbox}\to{\bblox}
\]
which is the identity on objects and it assigns to a morphism $f\colon\{0,1\}^n\to\{0,1\}^m$
the linear morphism $g\colon\{0,1\}^n\to\{0,1\}^m$ from Lemma~\ref{lem:linearization-auto-bblox} which is a functorial assignment since it is unique
and $\phi$ commutes as in Lemma~\ref{lem:linearization-auto-bblox}.
Clearly, $H$ exhibits the inclusion $j\colon{\bblox}\to{\bbox}$ as a retraction.

Now we can describe the functors which establish the equivalence $\int_{\bbGamma} \ast/\F_2^\bullet\simeq{\bbox}$ over $\bbGamma$ which will make for a commutative diagram
\begin{center}
\begin{tikzcd}
	\int_{\bbGamma} \ast/\F_2^\bullet & {{\bbox}} \\
	{\bbGamma} & {\bblox}
	\arrow["U", shift left, from=1-1, to=1-2]
	\arrow["{\mathrm{Un}}"', from=1-1, to=2-1]
	\arrow["V", shift left, from=1-2, to=1-1]
	\arrow["H", from=1-2, to=2-2]
	\arrow["G", shift left, from=2-1, to=2-2]
	\arrow["F", shift left, from=2-2, to=2-1]
\end{tikzcd}.
\end{center}
The functor $V$ acts as the identity on objects and it assigns to a morphism $f\colon\{0,1\}^n\to\{0,1\}^m$
\[
    V(f) = (F(H(f)),f(0)).
\]
For functoriality all that remains to be shown is that
\[
    f(g(0)) = f(0) + (F(H(f)))_\ast(g(0))
\]
in $\F_2^m$.\footnote{Which basically is affine linearity of $f$.}
But $F(H(f))_\ast = G(F(H(f))) = H(f)$ and so
\[
    (F(H(f)))_\ast(g(0)) = H(f)(g(0)) = \phi_{f(0)}(f(g(0)))
\]
for the unique $\phi_{f(0)}$ from Lemma~\ref{lem:linearization-auto-bblox}.
But in $\F_2^m$, the map $\phi_{f(0)}$ acts as addition of $f(0)$ and so
\[
    f(0)+\phi_{f(0)}(f(g(0))) = f(0) + f(0) + f(g(0)) = f(g(0)).
\]
It then automatically fulfills $\mathrm{Un}\circ V = F\circ H$.

The functor $U$ acts on objects identically and to a morphism $(f,a)$ it assigns $\phi_a\circ G(f)$.
Now functoriality follows from the statements made in Lemma~\ref{lem:linearization-auto-bblox}:
\begin{align*}
        U((f,a)\circ(g,b))
    &=  U(f\circ g,a+f_\ast(b)) \\
    &=  \phi_{a+f_\ast(b)}\circ G(f\circ g) \\
    &=  \phi_a\phi_{G(f)(b)}G(f)G(g) \\
    &=  \phi_a G(f)\phi_b G(g) \\
    &=  U(f,a)\circ U(g,b).
\end{align*}
In particular, $H\circ U = G\circ\mathrm{Un}$.
It is then easily verified that $U$ and $V$ form an equivalence of categories.
\end{proof}

\begin{remark}\label{rem:bbox-factorization}
In $\Fin_\ast$, every morphism $f\colon\langle m\rangle\to\langle n\rangle$
can be decomposed as
\begin{center}
\begin{tikzcd}
	{\langle m\rangle} & {\langle k\rangle} & {\langle n\rangle}
	\arrow["\alpha", from=1-1, to=1-2]
	\arrow["\sigma", from=1-2, to=1-3]
\end{tikzcd},
\end{center}
where $\alpha$ is the inert map that is the retraction of the order-preserving/canonical    
inclusion $\langle k\rangle^\circ\sub\langle m\rangle^\circ$ of the points $i$ where $f(i)\neq\ast$
and $\sigma$ is active.
This yields a decomposition of every map $f\colon\{0,1\}^n\to\{0,1\}^m$ in $\bbox$ as
\begin{center}
\begin{tikzcd}
	{\{0,1\}^n} & {\{0,1\}^k} & {\{0,1\}^m} & {\{0,1\}^m}
	\arrow["{\sigma^\ast}", from=1-1, to=1-2]
	\arrow["{\alpha^\ast}", from=1-2, to=1-3]
	\arrow["\phi", from=1-3, to=1-4]
\end{tikzcd}
\end{center}
where $\sigma^\ast$ is a (not necessarily bijective) permutation of coordinates arising from the active map $\sigma$,
$\alpha^\ast$ is a projection that only inserts zeros (arising from the inert $\alpha$)
and $\phi$ is the flip corresponding to a unique element $a\in\F_2^m$, in this case $a=f(0)$.
\end{remark}


\begin{example}
Let us draw the objects $n=1,2,3$ in the above categories.
Depicted are the lower-dimensional subobjects with their corresponding intersections.


\begin{center}
\pgfdeclarelayer{facefill}
\pgfdeclarelayer{afffacefill}
\pgfdeclarelayer{faceoutline}
\pgfdeclarelayer{afffaceoutline}

\pgfsetlayers{facefill,afffacefill,faceoutline,afffaceoutline,main}

\newcommand{\Face}[1]{%
  \begin{pgfonlayer}{facefill}
    \begin{scope}[blend mode=multiply]
      \path[fill=gray, fill opacity=0.12] #1 -- cycle;
    \end{scope}
  \end{pgfonlayer}
  \begin{pgfonlayer}{faceoutline}
    \path[draw=gray!60,line width=0.45pt] #1 -- cycle;
  \end{pgfonlayer}
}

\newcommand{\CornerFace}[1]{%
  \begin{pgfonlayer}{facefill}
    \begin{scope}[blend mode=multiply]
      \path[fill=gray, fill opacity=0.16] #1 -- cycle;
    \end{scope}
  \end{pgfonlayer}
  \begin{pgfonlayer}{faceoutline}
    \path[draw=gray!60,line width=0.45pt] #1 -- cycle;
  \end{pgfonlayer}
}



\newcommand{\AffFace}[1]{%
  \begin{pgfonlayer}{afffacefill}
    \path[fill=gray, fill opacity=0.10] #1 -- cycle;
  \end{pgfonlayer}
  \begin{pgfonlayer}{afffaceoutline}
    \path[draw=gray!60,line width=0.45pt] #1 -- cycle;
  \end{pgfonlayer}
}

\newcommand{\DiagFace}[1]{%
  \begin{pgfonlayer}{afffacefill}
    \path[fill=gray, fill opacity=0.16] #1 -- cycle;
  \end{pgfonlayer}
  \begin{pgfonlayer}{afffaceoutline}
    \path[draw=gray!70,line width=0.5pt] #1 -- cycle;
  \end{pgfonlayer}
}


\begin{tikzpicture}[
    x=1cm,y=1cm,
    font=\small,
    line cap=round,
    line join=round,
    vertex/.style={circle,fill,inner sep=1.4pt},
    ambient/.style={draw=gray!55, line width=0.45pt},
    mainedge/.style={draw=black, line width=0.8pt},
    affdiag/.style={draw=black, line width=0.65pt},
    simpedge/.style={
        draw=black,
        line width=0.75pt,
        -{Stealth[length=1.6mm,width=1.1mm]},
        shorten <= 8.0pt,
        shorten >= 8.0pt
    },
    rowlab/.style={anchor=east},
    collab/.style={font=\normalsize},
scale=0.6
]


\node[rowlab] at (-0.8,  0.0) {$\bbox$};
\node[rowlab] at (-0.8, -3.6) {$\bblox$};
\node[rowlab] at (-0.8, -7.4) {$\inert$};


\begin{scope}[shift={(0,0.4)}]
    \coordinate (A) at (0,0);
    \coordinate (B) at (2.2,0);

    \draw[mainedge] (A)--(B);

    \node[vertex] at (A) {};
    \node[vertex] at (B) {};
\end{scope}


\begin{scope}[shift={(4.3,0.0)}]
    \coordinate (A) at (0,0);
    \coordinate (B) at (2.5,0);
    \coordinate (C) at (0,2.5);
    \coordinate (D) at (2.5,2.5);

    \AffFace{(A)--(B)--(D)--(C)}

    \draw[mainedge] (A)--(B);
    \draw[mainedge] (B)--(D);
    \draw[mainedge] (D)--(C);
    \draw[mainedge] (C)--(A);

    \draw[affdiag] (A)--(D);
    \draw[affdiag] (B)--(C);

    \foreach \P in {A,B,C,D} {
        \node[vertex] at (\P) {};
    }
\end{scope}


\begin{scope}[shift={(9.25,0.0)}]
    \coordinate (A) at (0,0);
    \coordinate (B) at (2.4,0);
    \coordinate (C) at (0,2.4);
    \coordinate (D) at (2.4,2.4);

    \coordinate (E) at (1.1,0.8);
    \coordinate (F) at (3.5,0.8);
    \coordinate (G) at (1.1,3.2);
    \coordinate (H) at (3.5,3.2);


    \AffFace{(A)--(B)--(D)--(C)}
    \AffFace{(E)--(F)--(H)--(G)}

    \AffFace{(A)--(B)--(F)--(E)}
    \AffFace{(C)--(D)--(H)--(G)}

    \AffFace{(A)--(C)--(G)--(E)}
    \AffFace{(B)--(D)--(H)--(F)}


    \DiagFace{(A)--(B)--(H)--(G)}
    \DiagFace{(A)--(C)--(H)--(F)}
    \DiagFace{(A)--(E)--(H)--(D)}

    \draw[mainedge] (A)--(B);
    \draw[mainedge] (B)--(D);
    \draw[mainedge] (D)--(C);
    \draw[mainedge] (C)--(A);

    \draw[mainedge] (E)--(F);
    \draw[mainedge] (F)--(H);
    \draw[mainedge] (H)--(G);
    \draw[mainedge] (G)--(E);

    \draw[mainedge] (A)--(E);
    \draw[mainedge] (B)--(F);
    \draw[mainedge] (C)--(G);
    \draw[mainedge] (D)--(H);


    \draw[affdiag] (A)--(D);
    \draw[affdiag] (B)--(C);

    \draw[affdiag] (E)--(H);
    \draw[affdiag] (F)--(G);

    \draw[affdiag] (A)--(F);
    \draw[affdiag] (B)--(E);

    \draw[affdiag] (C)--(H);
    \draw[affdiag] (D)--(G);

    \draw[affdiag] (A)--(G);
    \draw[affdiag] (C)--(E);

    \draw[affdiag] (B)--(H);
    \draw[affdiag] (D)--(F);

    \draw[affdiag] (A)--(H);
    \draw[affdiag] (B)--(G);
    \draw[affdiag] (C)--(F);
    \draw[affdiag] (D)--(E);

    \foreach \P in {A,B,C,D,E,F,G,H} {
        \node[vertex] at (\P) {};
    }
\end{scope}

\begin{scope}[shift={(0,-3.6)}]
    \coordinate (A) at (0,0);
    \coordinate (B) at (2.2,0);
    \draw[mainedge] (A)--(B);
    \node[vertex] at (A) {};
    \node[vertex] at (B) {};
\end{scope}

\begin{scope}[shift={(4.3,-4.0)}]
    \coordinate (A) at (0,0);
    \coordinate (B) at (2.5,0);
    \coordinate (C) at (0,2.5);
    \coordinate (D) at (2.5,2.5);

    \Face{(A)--(B)--(D)--(C)}

    \draw[mainedge] (A)--(B);
    \draw[mainedge] (A)--(C);
    \draw[mainedge] (A)--(D);

    \foreach \P in {A,B,C,D} {
        \node[vertex] at (\P) {};
    }
\end{scope}

\begin{scope}[shift={(9.25,-4.4)}]
    \coordinate (A) at (0,0);
    \coordinate (B) at (2.4,0);
    \coordinate (C) at (0,2.4);
    \coordinate (D) at (2.4,2.4);

    \coordinate (E) at (1.1,0.8);
    \coordinate (F) at (3.5,0.8);
    \coordinate (G) at (1.1,3.2);
    \coordinate (H) at (3.5,3.2);

    \Face{(A)--(B)--(D)--(C)}
    \Face{(A)--(E)--(G)--(C)}
    \Face{(A)--(B)--(F)--(E)}
    \Face{(A)--(D)--(H)--(E)}
    \Face{(A)--(B)--(H)--(G)}
    \Face{(A)--(C)--(H)--(F)}

    \draw[mainedge] (A)--(B);
    \draw[mainedge] (A)--(C);
    \draw[mainedge] (A)--(D);
    \draw[mainedge] (A)--(E);
    \draw[mainedge] (A)--(F);
    \draw[mainedge] (A)--(G);
    \draw[mainedge] (A)--(H);

    \foreach \P in {A,B,C,D,E,F,G,H} {
        \node[vertex] at (\P) {};
    }
\end{scope}


\begin{scope}[shift={(0,-7.4)}]
    \coordinate (A) at (0,0);
    \coordinate (B) at (2.2,0);

    \draw[ambient] (A)--(B);
    \node[vertex] at (A) {};
    \node[vertex,fill=gray!45] at (B) {};
\end{scope}

\begin{scope}[shift={(4.3,-7.8)}]
    \coordinate (A) at (0,0);
    \coordinate (B) at (2.5,0);
    \coordinate (C) at (0,2.5);
    \coordinate (D) at (2.5,2.5);

    \draw[ambient] (A)--(B);
    \draw[ambient] (A)--(C);
    \draw[ambient] (A)--(D);

    \draw[mainedge] (A)--(B);
    \draw[mainedge] (A)--(C);
\CornerFace{(A)--(B)--(D)--(C)};

    \foreach \P in {B,C,D} {
        \node[vertex,fill=gray!45] at (\P) {};
    }
            \node[vertex] at (A) {};

\end{scope}

\begin{scope}[shift={(9.25,-8.4)}]
    \coordinate (A) at (0,0);
    \coordinate (B) at (2.4,0);
    \coordinate (C) at (0,2.4);
    \coordinate (D) at (2.4,2.4);

    \coordinate (E) at (1.1,0.8);
    \coordinate (F) at (3.5,0.8);
    \coordinate (G) at (1.1,3.2);
    \coordinate (H) at (3.5,3.2);

    \draw[ambient] (A)--(B);
    \draw[ambient] (A)--(C);
    \draw[ambient] (A)--(D);
    \draw[ambient] (A)--(E);
    \draw[ambient] (A)--(F);
    \draw[ambient] (A)--(G);
    \draw[ambient] (A)--(H);

    \CornerFace{(A)--(B)--(D)--(C)}
    \CornerFace{(A)--(B)--(F)--(E)}
    \CornerFace{(A)--(C)--(G)--(E)}

    \draw[mainedge] (A)--(B);
    \draw[mainedge] (A)--(E);
    \draw[mainedge] (A)--(C);

    \foreach \P in {A} {
        \node[vertex] at (\P) {};
    }
    \foreach \P in {B,C,D,E,F,G,H} {
    \node[vertex,fill=gray!45] at (\P) {};
    }
\end{scope}
\end{tikzpicture}
    \end{center}
\end{example}

From now on, we will freely switch between $\bbox$ and $\int_{\bbGamma} \ast/\F_2^\bullet$ depending on which category is better suited for certain applications.
Moreover, we will denote objects of $\bbox$ by $\langle n\rangle$, $\{0,1\}^n$ or ${\bbox}^n$,
depending on the mood of the authors when writing the corresponding parts.%
\footnote{As rough guidance: $\langle n\rangle$ is mostly used for emphasizing the abstract object,
and for the relation to $\bbGamma$,
$\{0,1\}^n$ is written for either emphasizing the relation to the geometric intuition of a cube
or for the combinatorial description of the category,
and $\bbox^n$ is mostly used as the Yoneda embedding of $\langle n\rangle$.}


\subsubsection{Properties of the Indexing Category}

Having seen their importance for nilspace theory,
let us recall some important properties of the categories $\Fin_\ast$, $\bbGamma$ and $\Fin_{\ast,\leq 1}$.

\begin{lemma}\label{lem:basic-prop-indexcat}
\begin{enumerate}[(i)]
    \item The category $\Fin_\ast$ is a Cisinski generalized Reedy category%
    \footnote{see \cite[Def 8.1.1]{Cisinski2006} or \cite{Berger2010}, and \cite{Riehl2014,Riehl2017} for the theory of Reedy categories}
        and has a symmetric monoidal structure given on objects by $\langle n\rangle\otimes\langle m\rangle = \langle n+m\rangle$.
    \item The category $\bbGamma$ is a Cisinski generalized Reedy category and has a symmetric monoidal structure given on objects by $\langle n\rangle\otimes\langle m\rangle = \langle n+m\rangle$.
    \item The category $\Fin_{\ast,\leq 1}$ is a Cisinski generalized Reedy category and has a symmetric monoidal structure given on objects by $\langle n\rangle\otimes\langle m\rangle = \langle n+m\rangle$.
\end{enumerate}
\noindent
Furthermore, in $\Fin_\ast$,
\begin{enumerate}[(1)]
    \item the epimorphisms agree with surjective maps and the split epimorphisms, and
    \item the monomorphisms agree with (everywhere defined) injective maps and the split monomorphisms.
\end{enumerate}
\end{lemma}

\begin{proof}
Let us begin with the latter assertions.
\begin{enumerate}[(1)]
\item If $\alpha\colon\langle n\rangle\to\langle m\rangle$ is an epimorphism in $\Fin_\ast$
one can test against $\langle m\rangle\to\langle 1\rangle$ to see that $\alpha$ is surjective,
and the existence of the section is evident.
\item If $\alpha\colon\langle n\rangle\to\langle m\rangle$ is a monomorphism in $\Fin_\ast$
one can test against $\langle 1\rangle\to\langle n\rangle$ to see that $\alpha$ is injective and everywhere defined.
Every injection has a retraction given by the identity on the image of $\alpha$ and undefined everywhere else.
\end{enumerate}
Now to the structures on the categories.
\begin{enumerate}[(i)]
\item The degree function is given by $\mathrm{d}(\langle n\rangle) = n$.
The degree increasing maps $\Fin_{\ast,+}$ are the monomorphisms
and the degree decreasing maps $\Fin_{\ast,-}$ are the surjections.
\item One simply takes $\bbGamma_+ = ((\Fin_\ast)_-)^\op$ and $\bbGamma_- = ((\Fin_\ast)_+)^\op$.
\item Here one uses the everywhere defined injections as $(\Fin_{\ast,\leq 1})_+$ and $(\Fin_{\ast,\leq 1})_- = \finert$ which correspond to the surjections in $\Fin_{\ast,\leq 1}$.\qedhere
\end{enumerate}
\end{proof}

\begin{remark}
Note that there does not exist a Cisinski generalized Reedy structure on $\finert$ with $\mathrm{d}(\langle n\rangle) = n$
since the only inert maps admitting a section are the isomorphisms, so $(\finert)_-$ is the core of $\finert$.
But simultaneously this implies that $(\finert)_+$ needs to consist of all maps.
This contradicts the fact that every inert map is surjective and thus lowers the degree.
The same argument works for $\inert$.
\end{remark}



\begin{lemma}\label{lem:fin-star-co-comp}
The category $\Fin_\ast$ is both finitely complete and finitely cocomplete
and the projection $\Fin_\ast\to\Fin$ creates finite limits and connected colimits.
\end{lemma}

\begin{proof}
This is well known.
Recall that the coproduct of finite pointed sets $A$ and $B$ is the pushout $A\sqcup_\ast B$.
\end{proof}

With this, we directly see that the monoidal structure on $\bbGamma$ is cartesian.

\begin{corollary}\label{rem:bblox-symm-monoid}
The symmetric monoidal structure on $\bblox$ defined by
\[
    \langle n\rangle \otimes\langle m\rangle =\langle n+m\rangle
\]
agrees with the cartesian one, i.e.,
\[
    \langle n+m\rangle=\langle n\rangle \times\langle m\rangle.
\]
\end{corollary}

Now, these properties translate through the unstraightening to the whole cube category $\bbox$ as follows.

\begin{lemma}\label{lem:bbox-epi-mono}
In the category $\bbox$, we have that
\begin{enumerate}[(i)]
    \item the epimorphisms agree with the surjective maps and with the split epimorphisms, and
    \item the monomorphisms agree with the injective maps and with the split monomorphisms.
\end{enumerate}
Furthermore, the monomorphisms are those maps in whose canonical representation
each $x_j$ appears (as $x_j$ or as $1 - x_j$) at least once.
The epimorphisms are precisely those maps such that no $y_i$ is constant and every $x_j$ appears at most once.
\end{lemma}

\begin{proof}
Let $f\colon{\bbox}^n\to{\bbox}^m$ be an epimorphism in ${\bbox}$.
We factor it as $\phi\circ g$ with $\phi\in\F_2^m$, and so we can assume that $g$ is in $\bbGamma$
and an epimorphism.
By Lemma~\ref{lem:basic-prop-indexcat}, its opposite $\alpha$ in $\Fin_\ast$ is an everywhere defined injection and there exists a splitting $\beta$ with $\beta\circ\alpha = \id$.
The splitting transfers to $f=\phi g$ since $\phi$ is invertible, and so $f$ is a split epimorphism.
Now if $f$ is a split epimorphism, it is in particular split in $\set$ and so is surjective.
On the other hand, every surjection clearly is an epimorphism (by applying the forgetful functor $\bbox\to\set$).
Moreover, from the canonical representation of $g = \alpha^\ast$ where $\alpha$ is an everywhere defined injection,
we infer that there appears no constant term $0$ and each $x_i$ appears at most once, since $\alpha$ has fibers of size at most one.
Therefore, the same holds for $f$.
On the other hand, every map $f$ with the same properties on its canonical representation is clearly surjective.

If $f\colon{\bbox}^n\to{\bbox}^m$ is a monomorphism in $\bbox$,
then in its factorization $f=\phi\circ g$, so is $g$.
Again, by Lemma~\ref{lem:basic-prop-indexcat}, we have that for $g=\alpha^\ast$ the map $\alpha$ is a split epimorphism, and hence $g$ is split monic.
The same line of reasoning as above also shows that split monomorphisms are injective,
and injections are monomorphisms.
Furthermore, $g=\alpha^\ast$ for $\alpha$ surjective shows that in the canonical representation of $g$,
each $x_i$ appears at least once.
And conversely, every such map is injective.
\end{proof}

\begin{lemma}\label{lem:bbox-reedy-cat}
The category $\bbox$ inherits a Cisinski generalized Reedy structure from $\bbGamma$
where the degree of $\{0,1\}^n$ is given by $n$,
the monomorphisms are the degree increasing maps and the epimorphisms are the degree decreasing maps.
\end{lemma}

\begin{proof}
Clearly, the non-invertible monomorphisms increase, and the non-invertible epimorphisms decrease the degree (since in $\Fin$, an injection (resp. surjection) between sets of the same cardinality is invertible).
Moreover, the isomorphisms preserve the degree and are both mono- and epimorphisms.
The factorization $f = m\circ e$ with $e$ epimorphic and $m$ monomorphic
follows from the corresponding factorization in $\bbGamma$ (Lemma~\ref{lem:basic-prop-indexcat})
which also implies uniqueness up to isomorphism.
Lastly, Lemma~\ref{lem:bbox-epi-mono} shows that every epimorphism has a section
and since the factorization $f = \phi\circ g$ is unique,
the set of sections determines the epimorphism uniquely by Lemma~\ref{lem:basic-prop-indexcat}.
\end{proof}

\begin{lemma}\label{lem:bblox-bbox-pres-limits}
The inclusion $j\colon\bbGamma\to\bbox$ preserves finite limits.
\end{lemma}

\begin{proof}
For products, we have seen that $\langle n\rangle\times\langle m\rangle = \langle n+m\rangle$.
Applying $j$ this corresponds to the set-theoretical product $\{0,1\}^n\times\{0,1\}^m = \{0,1\}^{n+m}$
with linear projections.
Given another $k$-cube with morphisms ${\bbox}^k\to{\bbox}^n$ and ${\bbox}^k\to{\bbox}^m$,
there exists a unique set-map $\{0,1\}^k\to\{0,1\}^{n+m}$ making the diagram
\begin{center}
\begin{tikzcd}
	& {\{0,1\}^{n+m}} & \\
	{\{0,1\}^n} && {\{0,1\}^m} \\
	& {\{0,1\}^k}
	\arrow[from=1-2, to=2-1]
	\arrow[from=1-2, to=2-3]
	\arrow[dashed, from=3-2, to=1-2]
	\arrow[from=3-2, to=2-1]
	\arrow[from=3-2, to=2-3]
\end{tikzcd}
\end{center}
commutative.
Looking at its $i$-th coordinate we see that, by commutativity of the diagram
and the fact that all solid arrows are in $\bbox$, it is of the form $0$, $1$, $x_{j(i)}$ or $1-x_{j(i)}$ for a unique $1\leq j(i)\leq k$.
This means that the set-map is in $\bbox$ and so $j$ preserves products.

For the preservation of equalizers consider the diagram
\begin{center}
\begin{tikzcd}
	{j\langle k\rangle} & {j\langle m\rangle} & {j\langle m'\rangle} \\
	{{\bbox}^n}
	\arrow["e", from=1-1, to=1-2]
	\arrow["f", shift left, from=1-2, to=1-3]
	\arrow["g"', shift right, from=1-2, to=1-3]
	\arrow["L"', from=2-1, to=1-2]
\end{tikzcd}
\end{center}
where $e$ is the equalizer of $f$ and $g$ in $\bbGamma$
and $fL = gL$, but $L$ is a map in $\bbox$.
Since $e$ is monic in $\bbGamma$, it is injective in $\set$.
Moreover, for all $a\in\{0,1\}^m$ with $f(a) = g(a)$ there exists $c\in\{0,1\}^k$ such that $e(c) = a$:
If this weren't the case, we could define a map $\langle 1\rangle\to\langle m\rangle$ in $\bbGamma$ mapping $1$ to $a$
which then would get equalized by $f$ and $g$, but couldn't factor through $e$, a contradiction.
Hence, upon applying the forgetful functor to $\set$,
this diagram is an equalizer diagram.
Hence there exists a unique map $T\colon\{0,1\}^n\to\{0,1\}^k$ of sets(!) making the diagram
\begin{center}
\begin{tikzcd}
	{j\langle k\rangle} & {j\langle m\rangle} & {j\langle m'\rangle} \\
	{{\bbox}^n}
	\arrow["e", from=1-1, to=1-2]
	\arrow["f", shift left, from=1-2, to=1-3]
	\arrow["g"', shift right, from=1-2, to=1-3]
	\arrow["T", dashed, from=2-1, to=1-1]
	\arrow["L"', from=2-1, to=1-2]
\end{tikzcd}
\end{center}
commutative (in $\set$).
Now since $e$ is injective, from Lemma~\ref{lem:bbox-epi-mono} we infer that the $i$-th coordinate of $T$ is of the form $0$, $1$, $x_{j(i)}$ or $1-x_{j(i)}$ for a unique $1\leq j(i)\leq n$.
This means that the set-map $T$ is in $\bbox$ and so $j$ preserves equalizers.
\end{proof}


\begin{corollary}\label{prop:bbox-symm-monoid}
There is a symmetric monoidal structure on $\bbox$ defined by the formula
\[
    {\bbox}^n\otimes{\bbox}^m={\bbox}^{n+m}
\]
which agrees with the cartesian monoidal structure.
\end{corollary}

\begin{proof}
All that is to be shown is that we have ${\bbox}^n\times{\bbox}^m\simeq{\bbox}^{n+m}$.
But this follows from Lemma~\ref{lem:bblox-bbox-pres-limits} and Lemma~\ref{rem:bblox-symm-monoid}.
\end{proof}

\begin{remark}
Note that the property that the symmetric monoidal structure is the cartesian one
is false in other classical indexing categories, e.g., in the simplex category $\bbDelta$
or other variants of cube categories in cubical homotopy theory, see below.
\end{remark}


\begin{remark}
Note that Segal's category $\bbGamma$ is used to model $\bbGamma$-spaces,
see \cite{Bousfield1978,Brown1981,Schwede1999,Schwede2000},
as well as being the (opposite of the) $\E_\infty$-operad modeling commutative monoids in $\infty$-categories,
see \cite{Lurie2017,Hebestreit2021}.
The trivial operad $\finert$ in turn is used to define operads ($\infty$-multicategories).

Often, $\bbGamma$ is understood as a commutative variant of the classical simplex category $\bbDelta$
(which even is an Eilenberg-Zilber category).
One can relate them using the \emph{cut functor} of Dedekind cuts $\bbDelta\to\bbGamma$.
This functor exhibits every symmetric operad as a non-commutative operad
and shows that $\mathbb{E}_\infty$-monoids are $\mathbb{E}_1$-monoids, see \cite{Lurie2017,Hebestreit2021}.
There are also relations to other important indexing categories,
e.g. the trees of dendroidal sets, see \cite{Cisinski2011,Heuts2022}.
\end{remark}

\begin{remark}\label{def:cubical_homotopy_cubes}
There already is a well-established cubical homotopy theory, see e.g. \cite{Brown1981,Cisinski2006,Brown2011,Doherty2024}.
The cube category used in the theory of nilspaces has many more morphisms (diagonals)
than these cube categories.
However, it misses the usual \enquote{connections}.

In particular, there are Cisinski model structures to model spaces, and Joyal model structures to model $(\infty,n)$-categories using cubical sets.
However, those structures do not fit our purposes, and in particular their definitions of fibration and homotopy group differ from ours.
\end{remark}

